% ======================================================================
\section{Hashing (for data indexing)}

% ======================================================================
\subsection{A subjective history of {\em ``Where did I put the damn keys?''}}

% ----------------------------------------------------------------------
\begin{frame}[fragile]{Searching, by hand}

  Earliest alternative to linear search: hand-built index
  \begin{itemize}
  \item Mesopotamia

    Clay tablets organized in collections, lexical catalogues constructed {\em systematically}.

  \item Ancien Egypt

    Administrative records by category, scribes produce {\em ordered} word (and sign) lists.

    \begin{itemize}
    \item Most often: single-level indexing
    \end{itemize}

  \item Greco-Roman world: most notable example is the Library of Alexandria
    \begin{itemize}
    \item Callimachus' {\em Pinakes}: huge catalogue by author {\em and} subject, with biblio info
    \item Likely the first known database
    \item ``Practical'' multi-level indexing
    \end{itemize}
  \item Romans credited to leverage alphabetical order for organization and search
    \begin{itemize}
    \item ``Abstract'' multi-level indexing
    \end{itemize}
  \item Pre-Google Internet had collaborative construction of curated website collections
    \begin{itemize}
    \item Still used on the Darknet today (darksites not discoverable by construction)
    \end{itemize}
  \end{itemize}

\end{frame}


% ----------------------------------------------------------------------
\begin{frame}[fragile]{Searching, with a machine}

  Main problems with an index: built by hand, logarithmic up to index depth.  

  You may still have a long list of data to check once you reached an index leaf.

  \begin{itemize}
  \item Mid-1940's: {\em Associative access} (making the goal explicit)

    K. Zuse (Germany), V. Bush (USA).
    
    {\em ``What if those strange machines allowed to {\em search data by its very contents}?''}

    \begin{itemize}
    \item Goal: automatic (non-human) scribes.
    \end{itemize}
  \item 1953: {\em Chaining hashtables} (closed addressing)

    H.P. Luhn.

    {\em ``Let's {\em compute from the data itself} the number of the list where it belongs!''}

    \begin{itemize}
    \item On an IBM 701: 2000 words of magnetic-drum memory, 72-bit machine words;
    \item Full pointers too costly: have smaller indices into storage array instead;
    \item That's basically our pointer arithmetics (to save on the pointers size.)
    \end{itemize}
  \end{itemize}

\end{frame}


% ----------------------------------------------------------------------
\begin{frame}[fragile]{Searching, with a machine of today}

  \begin{itemize}
  \item 1954: {\em Linear hashtables} (open addressing)

    G. Amdahl, E. McGraw, A. Samuel (IBM 701) + A. Ershov independently.

    {\em ``By the way, you could even have {\em implicit} lists!''}

  \item Late 1960's: The word ``hashing'' eventually becomes standard terminology
    \begin{itemize}
    \item Many variants explored: quadratic hashing, double-hashing, Cuckoo hashing, {\em etc.}
    \end{itemize}

  \item 1986: {\em Robin-Hood hashing} (minimize inequalities to maximize robustness)

    P. Celis.
    
    {\em ``Let's have a linear table {\em self-organize to maintain highest performance}!''}

    \begin{itemize}
    \item Robustness? {\em RH maintains highest performance under a larger panel of workloads.}
    \end{itemize}
    
  \item 2022: {\em SwissTable} (leverage insn-level, parallel probing)

    Abseil (company).
    
    {\em ``Let's have {\em ultra-fast SIMD lookups with metadata} control bytes!''}

    \begin{itemize}
    \item Robin-Hood on steroids (most lookups with a single probing insn)
    \item Influential design (Rust's HashBrown table, {\em etc.})
    \end{itemize}
  \end{itemize}
  
\end{frame}


% ======================================================================
\subsection{A world of complexity}

% ----------------------------------------------------------------------
\begin{frame}[fragile]{The big picture (searching)}

  Applications fall in two categories:
  \begin{enumerate}
  \item If you are only concerned with {\em exact search}, go with hashtables

    {\em ``What is the value of this key?''}

    Hashtables are unbeatable at this game (amortized constant-time)
    
  \item If you need {\em interval search}, go with ordered search trees

    {\em ``What are the keys in this interval?''} (spatial data, ecological data, {\em etc.})

    \begin{itemize}
    \item {\em Lower-bound} search primitive.
    \item The lower bound in the data may or may not be that of the search interval
    \item Easy to search for absence or presence+value of one specific key

      To instantiate exact search for comparison purposes.
    \end{itemize}
  \end{enumerate}

  That's a lot of ground to cover!
  
\end{frame}


% ----------------------------------------------------------------------
\begin{frame}[fragile]{The big picture (hashing)}

  We are {\em not} doing crypto!
  \begin{itemize}
  \item Hash data to produce secure digests (Keccak, {\em etc.})
  \item Much more stringent mathematical requirements

    Collision resistance, pre-image resistance, second pre-image resistance.
  \end{itemize}

  But we may still have adversaries:
  \begin{itemize}
  \item DoS (Denial-of-Service) attack on routers

  \item Targeted weakness was the router internal hashtable

  \item Hackers triggering catastrophic, degenerate linear-time search

    Poorly designed hash function allowed crafting the right attack sequence

  \item SipHash hash function specifically designed as a counter-measure
  \end{itemize}

\end{frame}


% ----------------------------------------------------------------------
\begin{frame}[fragile]{The big picture (hashing + searching)}

  \begin{itemize}
  \item A hash is a source of pseudo-randomness derived from the key alone.
  \end{itemize}

  Search trees may use a hash in essentially two forms:

  \begin{enumerate}
  \item Cheap and efficient randomized balancing of a BST (treap)

    More complicated and stringent balancing with AVL trees, Red-Black trees.
  \item Plain routing info inside the tree (hashtrie)

    Balancing unordered keys.
  \end{enumerate}

  It is common practice for a hash-based search structure to {\em finalize} the hash.

  The user must provide the hash function, which may be poorly designed:
  \begin{itemize}
  \item Add some internal layer of ``hashing the hash''
  \item Must be fast (a few CPU insns): \textsc{xor}-based bit-mixer, Fibonacci hashing
  \end{itemize}

  Cheap patch against entry-level contingencies, close to useless under attack.
  \begin{itemize}
  \item Gain operational reach and stability.
  \end{itemize}
  
\end{frame}


